\section{The measure}\label{sec:measure}\label{sec:construct}

\subsection{Two readers, one living backwards}

Imagine two readers of a filing's goods/services description, each having seen
a disjoint half of its industry. One has read only what the industry filed in
the five years before it. The other, living backwards as Merlin was said to,
has read only the five years after. Each is asked how surprising the wording
is. The measurement is how much they disagree: a description that startles the
first reader and not the second arrived before its industry's language moved
that way.

Both windows are anchored on the filing's own date, spanning the 1,826
days before it
and the 1,826 days after; both exclude the filing date itself, so a filing
never appears in the reference of another filing made the same day, and no
two filings made on different days share a reference. Inside a window every day counts the same. The measure asks whether
wording is unusual against a period, not whether it is ahead of a trend
inside that period; it has no momentum channel.

\subsection{What is measured}

The quantity is information content in Shannon's sense. Something that
happens with probability $p$ carries $\log(1/p)$ of information when it
happens: a common thing carries little, a rare thing a lot. Here the things
are the themes a description draws on, and the probabilities are how often
the filing's own industry drew on each of them in the reference window. A
filing's score asks, in effect: how surprising is this mix of themes to
someone who has only seen what the industry files? Formally it is the
average information content of the filing's themes under its industry's
frequencies, weighted by how heavily the filing uses each, less the same
average taken under the filing's own frequencies.

The apparatus is that of \citet{Murdock2017}, developed on Darwin's reading
notebooks, and \citet{Barron2018}, who extended it to French Revolution
parliamentary debates. They score a document against those before it and
again against those after, calling the two quantities \emph{novelty} and
\emph{transience} and their signed difference \emph{resonance}. Their object
was the life of ideas: how often an idea appeared, and how ideas stood in
relation to one another, across one reading life or one assembly's debate.
The object here is the composition of commercial offerings, which components
a product or service is described as having, and how that composition shifts
across an industry over time. The arithmetic carries over unchanged. The
interpretation does not: their document has one author and the reference
stream is that author's own reading, so a high value says something about a
mind; here a filing is scored against other firms' filings, and its author is
usually counsel drafting for scope of protection. The same shift renames
what the model finds. In those settings, and in the name of the method,
latent Dirichlet allocation identifies \emph{topics} --- what a document
is about. A goods/services description is not about anything in that
sense; it lists components, and across millions of filers a cluster of
terms that co-occur (\emph{video, music, games, electronic, audio}) is a
component many offerings share --- a theme the entrants' offerings
evidence, not a subject. \emph{Theme} is the word used throughout;
\emph{topic} appears only in the model's name.

\subsection{From a description to two numbers}\label{sec:definitions}

\paragraph{Counting words.} A description is reduced to a bag of words: the
text is lowercased, split into tokens of three or more letters, and counted,
with adjacent pairs counted as tokens in their own right so that
\emph{computer software} registers alongside \emph{computer} and
\emph{software}. Order is discarded. Take the filing for \textsc{Amazon.com},
lodged in the advertising and retail class on 18 April 1997:

\begin{quote}\small
computerized on line search and ordering service featuring the wholesale and
retail distribution of books, music, motion pictures, multimedia products and
computer software in the form of printed books, audiocassettes,
videocassettes, compact disks, floppy disks, CD ROMs, and direct digital
transmission
\end{quote}

\noindent Of its 45 distinct terms, 36 are in the theme model's vocabulary,
among them \emph{computerized}, \emph{search}, \emph{ordering},
\emph{ordering service}, \emph{retail}, \emph{distribution}, \emph{printed
books} and \emph{digital transmission}.

\paragraph{Grouping words into themes.} Each filing is then re-described
in a fixed, much smaller set of coordinates. We run latent Dirichlet
allocation once over a stratified sample of the whole corpus, all 45
classes together. The model develops a list of 50 themes --- clusters of
words that tend to occur together --- and infers, for any filing, to what
degree its description invokes each one. Nothing supervises this; the
clusters come from co-occurrence alone, and a theme is read off from its
leading words: one collects the vocabulary of media and entertainment
goods, another retail store services, another downloadable and online
software. Re-described this way, a filing becomes a mix of themes, and
two kinds of question open up: how a theme fares over time within a class
--- how much of the class's language an online-offering theme accounts
for, year by year --- and how the filings that invoke it fare. This
build --- fifty global themes, per-filing five-year windows, flat
weighting --- is the production scoring, used for every estimate unless
the text states another. The themes are shared across classes; what is
specific to a class is the reference: the average theme proportions of
that class's own filings in the window. The Amazon filing draws on six
themes at more than a twentieth of its weight: $0.41$ on media and
entertainment goods (\emph{video, computer, electronic, music, games,
audio}), $0.15$ on computer and information services (\emph{services,
computer, software, information, data}), $0.13$ on retail store services
(\emph{store, retail, store services, featuring}), $0.10$ on printed matter
(\emph{paper, books, printed, cards, stationery}), $0.09$ on downloadable and
online software, and $0.07$ on marketing of consumer goods. The online
appendix lists every theme with its leading words and its weight in each
class.\footnote{\onlineappendixnote}

Fifty is coarse, roughly one theme per industry.
Section~\ref{sec:three_scorings} refits the model at 500 themes and
separately at 50 themes within each class and repeats the five-year-proof
analysis under each; Section~\ref{sec:robust} reports the resolution sweep
and a second fit at the same resolution with a different seed. Fifty is
kept elsewhere because it is the coarsest resolution at which every result
holds, and because its software-class lead contrast sits at the low end of
the resolutions tried (only the 100-theme refit returns less).

\paragraph{Comparing against the two references.} Write $P$ for the filing's
distribution over themes, and $Q^{-}$ and $Q^{+}$
for that distribution averaged over every same-class filing in the two
windows, $[d_i - W,\, d_i)$ and $(d_i,\, d_i + W]$, with $W = 1{,}826$ days.
How far $P$ sits from a reference $Q$ is the Kullback--Leibler divergence,
\[
  \mathrm{KL}(P \,\|\, Q) \;=\; \sum_{k} P_k \log \frac{P_k}{Q_k},
\]
a sum over the 50 themes, zero when the two distributions agree and growing as
the filing puts weight where the industry puts little. It is measured in nats,
the natural-log unit. Write
\[
  K^{-} = \mathrm{KL}(P \,\|\, Q^{-}),
  \qquad
  K^{+} = \mathrm{KL}(P \,\|\, Q^{+}),
\]
so $K^{-}$ is what the first reader feels and $K^{+}$ the second: the
filing's past-facing and future-facing surprise.

\paragraph{Taking the difference.} Amazon's filing is scored against 41{,}166
earlier filings and 146{,}532 later ones, giving $K^{-} = 1.286$ and
$K^{+} = 1.163$. The two axes used throughout are their average and their
signed difference:
\[
  \underbrace{A \;=\; \tfrac{1}{2}\left(K^{-} + K^{+}\right)}_{\text{atypicality}},
  \qquad
  \underbrace{L \;=\; K^{-} - K^{+}}_{\text{lead}} .
\]
For Amazon, $A = 1.225$ and $L = +0.123$: against its own class, its
atypicality is at the 47th percentile and its lead at the 94th. The
industry moved toward this description in the five years after it was filed
and had not been writing that way in the five years before.
Table~\ref{tab:nomenclature} maps the terms and their nearest occupants
in the literature; the signed difference is written $L$ throughout and
its magnitude $|L|$.

What ``early'' means is visible in the themes themselves: the three themes
that carry most of Amazon's description --- retail store services, media
and entertainment goods, and the still-small online-software theme --- all
became more common in class 035 over the following five years, the last
nearly doubling, while the two minor themes it touched receded. The filing
is closer to the class's future composition than to its past one; the
supplementary material tabulates the six theme shares.

Why the average and the difference, rather than the two surprises
themselves? Because $K^{-}$ and $K^{+}$ correlate at $0.988$ on the
software class: a filing far from its class's past is almost always far
from its future too, so entering both is close to entering one variable
twice. The average collects that large shared component --- how far from
the class the filing sits --- and the difference collects the small
residual that carries the timing. Nothing more than a sum and a
difference is taken, with no weights estimated, and the two resulting
regressors are nearly uncorrelated ($+0.036$ against each other, where
$K^{-}$ alone would give $+0.114$). The cost is that $L$ is a small residual of
two large quantities, about a sixth of either level's spread, which
Section~\ref{sec:robust} quantifies.


\paragraph{Why the words are not scored directly.} The themes are an extra
step, and the words themselves could be compared to the reference: build the
filing's distribution over the class vocabulary from the words it uses, do
the same for every filing in the window, and take the divergence. The paper
reports that scoring as a cross-check. It cannot be the primary one, because
a divergence measured on words mostly counts the length of the filing. The
divergence is an average surprise less a spread: for any two distributions,
\[
  \mathrm{KL}(P \,\|\, Q) \;=\; \underbrace{-\textstyle\sum_k P_k \log Q_k}_{\text{how unexpected these words are}}
  \;-\; \underbrace{\left(-\textstyle\sum_k P_k \log P_k\right)}_{\text{how spread out the filing is}} ,
\]
and the spread of a filing over words is bounded by how many distinct words it
has: a description using $k$ terms cannot spread over more than $k$ of them,
so its spread cannot exceed $\log k$, while the reference is spread over about
900 terms. A five-word description is charged the whole difference in width,
roughly five nats, before any of its content is considered. To show this
directly, we took 3{,}000 software-class filings of at least 150 words,
scored each against one fixed reference, then cut each down to a random
handful of words and rescored it. Across a fiftyfold cut in length, the
surprise term moves by $0.016$ nats while the spread term moves by
$2.9$: essentially all of the change in the score is the length, not the
content. Synthetic filings drawn at random from the reference itself
track $\log(\text{distinct terms})$ to within a fiftieth of a nat
(supplementary material). Term-scored atypicality correlates with
the log of a filing's distinct term count at $-0.68$ in the software
class.

Themes break the coupling because the 50 coordinates are always occupied.
Under term scoring the effective number of coordinates a filing spreads over
rises seventeenfold from the shortest descriptions to the longest; under
themes it varies by less than a third, and not monotonically: a three-word
fragment gives the model little evidence, so its inferred proportions stay
close to the prior and spread across many themes. The spread term no longer
tracks length, and the score is dominated by the surprise term.
Theme-scored atypicality correlates with the log term count at $-0.22$ in
the software class and $-0.10$ in advertising and retail, against $-0.68$
and $-0.62$ under terms.

\paragraph{What the themes see, and what they do not.} The score compares a
filing's theme proportions to the average proportions of its industry, so it
reads which themes a filing draws on and how heavily. It does not, in any
material way, read how a filing combines them. To check this, we took
each filing's two leading themes on a 150{,}000-filing software-class
sample and asked how often that pair occurs together, relative to how
often the two themes occur separately. How rare the themes are
individually explains 70\% of the variance in atypicality; how unusual
the pairing is adds under two points on top of it. A filing assembled from ordinary parts in a
genuinely new configuration is therefore close to invisible to this measure.

Combinations can be scored in their own right, at two levels
(construction detail in the supplementary material). The first is
\emph{pair lead}: whether the class later moved toward the filing's
arrangements of themes. In the software class it is penalised much as
lead is ($+2.3$ points at the five-year proof, $t = 8.4$, while
correlating with lead at only $-0.16$). But it does not generalize.
Corpus-wide it is zero raw and $+0.5$pp net of lead and atypicality
(owner-clustered $t = 4.2$), significantly positive in 13 of 43 classes
and significantly negative in 3 --- including $-2.1$ in advertising and
retail --- so we report it as a class-specific bound rather than a second
construct. The second is word-level \emph{recombination}: adjacent word
pairs that are new to the class while both words are established there,
novelty assembled from familiar components
\citep{Fleming2001, Uzzi2013}. In the same class it is rewarded: the top
fifth fails $2.8$ points less with lead and atypicality held
($t = -10.4$), and it is uncorrelated with pair lead ($r = -0.01$). The
construct sees neither arrangement channel, and where they act they run
in opposite directions, so its blindness to combination is a bounded
omission rather than a hidden bias in one direction.

\subsection{Within class, not across it}

Goods/services descriptions are drafted to secure scope of protection, not to
describe a product distinctively, and counsel reuse language from the USPTO's
own Acceptable Identification of Goods and Services Manual, from house
precedent, and from competitors' successful filings. Identical wording across
unrelated firms is therefore common, and it carries through to the scores:
17\% of scored registrations issued 2002--2018 share a value with another
filing in the same class and registration year. Within a class, the measure picks up departure from a shared
drafting convention. Across classes it does not: a class where the Manual
supplies dense standard language shows a compressed distribution of
atypicality, and one where firms describe novel services in prose a wide one,
for reasons unrelated to how innovative either industry is. Every estimate is
therefore taken within class and year, and magnitudes compare across classes
only as signs and orderings.

The vocabulary is positional rather than causal: a leading filing was
ahead of its industry's language, which implies a trajectory without
asserting the filing caused it.

\subsection{Term scoring measures length; theme scoring does not}\label{app:representation}

Figure~\ref{fig:representation} puts the two representations side by side
over the same axes, with each hexagon coloured by the mean length of the
filings in it. Under term scoring the colour runs from dark near the
origin --- a median of about 200 terms --- to near-white at the far end,
where the median is eight. That gradient is the measurement problem: a
filing scores as atypical largely because it is short. Under theme scoring
the same surface is a narrow band of roughly uniform tone.

Two filings are marked in both columns. Under the production scoring they
behave differently from one another, which is the point. \textsc{Amazon.com}
(1997) is at the 98th percentile of lead on terms and the 94th on themes,
but only the 40th and 47th
percentiles of atypicality, so both agree that it was early and neither
mistakes it for strange. Apple's \textsc{iPod} (2001) divides them: the 31st
percentile of lead on terms against the 62nd on themes. The surfaces plotted
here are on the class-year build the figure was generated from, where the
length gradient is starkest; the percentiles just quoted are on the per-filing
build used for every estimate.

\begin{figure}[!htbp]\centering
\includegraphics[width=\textwidth]{results/representation_appendix.png}
\caption{Mean filing length across the two axes, by representation and
class, for filings 1995--2019 scored on per-filing reference windows. The
horizontal axis is surprise against the preceding five years, the vertical
axis surprise against the following five; a filing on the dashed diagonal
is equally far from both, and distance below it is lead. Hexagons holding fewer than 25 filings are blank; the retained
hexagons cover over 99.4\% of filings in every panel. The colour scale is
shared across panels; the axis ranges are not, because term-scored and
theme-scored KL are not on a common scale, so only shape and colour
gradient are comparable across columns. Both
columns are drawn on the filings the per-filing windows score, which
differ between representations by under one percent.}
\label{fig:representation}
\end{figure}

Each representation has the virtue the other lacks: theme scoring is
independent of document length, which term scoring is not, while term
scoring resolves novel word \emph{combinations}, which theme scoring does
not. We therefore use theme scoring for every estimate, because the
length confound contaminates firm-level comparisons, and term scoring for
every worked example, because the theme basis has no coordinate for the
language the examples turn on.

\subsection{How three filings encode}\label{app:exhibit}

Table~\ref{tab:exhibit} traces three filings through term scoring and theme
scoring at $T \in \{50, 200, 500\}$ --- three kinds of novelty, and how each
representation handles them.

\textsc{Amazon.com} (1997, advertising/retail) is \emph{combinatorial}
novelty: every word is ordinary, but the bigrams assemble an invention
(``computerized line,'' ``line search,'' ``search ordering''; class
document frequencies 379, 134, 355). The filing's text reads ``on line''
as two words, the standard spelling in 1997, before \emph{online}
lexicalized into a single token; the analyzer drops the stopword ``on''
and the bigram ``line search'' survives
stopword removal. The novelty was necessarily combinatorial
because the word for it did not yet exist. Term scoring sees these
bigrams; under theme scoring at every resolution all of them are out of
vocabulary and the filing reads as a media retailer that uses computers.

\textsc{Kombucha} (1997, beer/soft drinks) is \emph{lexical} novelty: the
new thing arrives as new words (\emph{kombucha} df 567, \emph{saccharose}
out of vocabulary even at the class level). The word \emph{kombucha} is in
the theme vocabulary but, with no kombucha-like theme, the filing loads on
a generic beverages theme at every $T$ (Table~\ref{tab:exhibit}).

\textsc{Ethereum} (2018, software) is \emph{wave} novelty: emergent
vocabulary already diffusing (\emph{blockchains}, \emph{distributed
computing}). It loads on generic software and services themes at every
resolution, including $T = 200$, which does contain a
cryptocurrency/blockchain theme; there, the eleven occurrences of
\emph{software} dominate the theme assignment.

\begin{table}[!htbp]\centering\small
\caption{Top theme loading for three filings, by representation. Term
scoring resolves the novel content of all three; theme scoring resolves
none of it at any tested resolution.}
\label{tab:exhibit}
\begin{tabular}{lp{3.1cm}p{3.1cm}p{3.1cm}}
\toprule
& \textsc{Amazon.com} & \textsc{Kombucha} & \textsc{Ethereum} \\
& (1997) & (1997) & (2018) \\
\midrule
novel content (term) & line search, search ordering & kombucha, saccharose & blockchains, distributed computing \\
$T=50$ & video, computer, music (0.40) & beverages, fruit, drinks (0.64) & software, online, downloadable (0.33) \\
$T=200$ & entertainment, video, games (0.29) & beverages, drinks, coffee (0.60) & field, services, development (0.37) \\
$T=500$ & computer, software, data (0.22) & beverages, drinks, fruit (0.61) & software, online, computer (0.49) \\
\bottomrule
\end{tabular}
\end{table}

The pattern across the row is the measurement point in miniature: theme
scoring cannot represent combinatorial novelty (Amazon), dilutes lexical
novelty (Kombucha), and majority-votes away wave novelty (Ethereum) even
when a fitting theme exists. That the mark-level outcome results are
nearly identical under both scorings despite this blindness is what
licenses theme scoring as a robustness instrument; that the firm-level
results diverge is what exposes them as artifacts of the lexical
specificity term scoring additionally captures
(Section~\ref{sec:robust_hazard}).

\subsection{Scores are interpretable from 1995}\label{sec:burnin}

\emph{Burn-in} is the span at each end of the corpus where the reference
windows cannot be filled, so scores there reflect a thin reference rather
than the text being scored.

The signal depends on having enough past and future filings to populate
the references. The 5-year past window is incomplete before 1990 and the
future window after 2020, and both zones are excluded
(Figure~\ref{fig:burnin}). The burn-in cut is set empirically. A thin
past reference inflates past-facing surprise, and profiling that bias
across all 44 sufficient-volume classes gives a median absolute bias of
$\sim$2.0 nats in 1990, 0.25 in 1991, 0.10 in 1992, and 0.04 from 1993
onward --- below any effect size interpreted in this paper. The
1995--1997 deviations ($\sim$0.15--0.17) are not burn-in: mechanical
contamination declines monotonically as the windows fill and cannot
rebound, so these reflect a genuine vocabulary shift in the corpus
during the dot-com period. Estimating a separate burn-in length per
class turns out to track each class's noise floor rather than its
reference density, so a single standard cut at 1995 is used,
conservative by about two years.

\begin{figure}[!htbp]\centering
\includegraphics[width=0.85\textwidth]{results/burnin.png}
\caption{Reference-window density by year on log scale. Hatched regions
are the past-side and future-side burn-in zones where the 5-year window is
incomplete. The dotted line marks the 1995 analysis cut. The windows
decide the score at the filing date; the gate outcome is read separately,
at registration age six to seven.}
\label{fig:burnin}
\end{figure}

